\documentclass[aps,prd,preprint,amsmath,amssymb]{revtex4-2}
\usepackage{graphicx}
\usepackage{hyperref}
\usepackage{booktabs}

\title{Vacuum Depletion as the Origin of the CKM Unitarity Deficit:\\
A Lattice QCD Perspective within the Aeternvm Vacuvm Framework}
\author{Gustavo Alves Conde}
\affiliation{Aeternvm Vacuvm Collaboration, Itagua\c{c}u, ES, Brazil}
\date{\today}

\begin{document}
\maketitle

\begin{abstract}
The recent Lattice QCD determinations of the $\gamma W$-box correction $\Delta_R^V$ have shifted the extracted value of $|V_{ud}|$, leading to a $\sim 3\sigma$ deficit in first-row CKM unitarity: $|V_{ud}|^2+|V_{us}|^2+|V_{ub}|^2=0.9985(5)$. We reinterpret this deficit within the Aeternvm Vacuvm framework of a depletable vacuum. We show that a vacuum-depletion coupling $\kappa_{\rm vac}=5\times10^{-4}$ with mean-field expectation $\langle\phi\rangle\sim 1.2$ naturally reproduces the observed deficit without invoking beyond-Standard-Model particles. A modified lattice action $S_{\rm Aet}=S_{\rm QCD}+\kappa_{\rm vac}\phi\,{\rm Tr}(F^2)$ predicts that the deficit survives the continuum limit $a\to0$, in contrast to ordinary discretization artifacts. Illustrative parameter scans and continuum extrapolations are presented.
\end{abstract}

\section{Introduction: The Cabibbo Angle Anomaly}
The Cabibbo Angle Anomaly (CAA) has become one of the most precise low-energy hints of physics beyond the Standard Model. Combining superallowed $0^+\to0^+$ nuclear $\beta$ decays~\cite{Hardy2020}, recent dispersive evaluations of $\Delta_R^V$~\cite{Seng2018,Seng2020}, and Lattice QCD calculations of the $\gamma W$ box~\cite{FermilabMILC}, one finds
\begin{equation}
\Delta_{\rm CKM}\equiv 1-\bigl(|V_{ud}|^2+|V_{us}|^2+|V_{ub}|^2\bigr)=0.0015(5).
\end{equation}
The conventional interpretation invokes right-handed currents, scalar operators, or other BSM contributions. In the Aeternvm Vacuvm framework we propose an alternative: the deficit reflects the energetic cost of deforming a depletable vacuum.

\section{Lattice Methodology and Aeternvm Extension}
Standard Lattice QCD extraction of the $\gamma W$-box proceeds through five steps:
\begin{enumerate}
\item Discretization on $L^3\times T$ lattices with spacings $a\simeq0.06$--$0.12\,{\rm fm}$;
\item Hybrid Monte Carlo sampling of the gauge ensemble weighted by $e^{-S}$;
\item Inversion of the Dirac operator;
\item Evaluation of the four-point correlator $\langle J_{\rm em}J_{\rm weak}\rangle$;
\item Continuum and infinite-volume extrapolations $a\to0$, $L\to\infty$.
\end{enumerate}

We introduce a minimal extension of the action,
\begin{equation}
S=S_{\rm QCD}+\kappa_{\rm vac}\sum_x\phi_x\,{\rm Tr}(F_{\mu\nu}F^{\mu\nu})+V(\phi),
\end{equation}
where $V(\phi)=\lambda(\phi^2-v^2)^2$ with $v=1$ and $\phi$ is a real scalar field representing the local vacuum density. The effective weak coupling is then
\begin{equation}
V_{ud}^{\rm eff}=V_{ud}^{\rm bare}\bigl(1-\kappa_{\rm vac}\langle\phi\rangle\bigr).
\end{equation}

\section{Illustrative Parameter Scan}
We performed a scan over 72 combinations of lattice spacing, volume, $\kappa_{\rm vac}$ and $\langle\phi\rangle$ (see ancillary file \texttt{varredura\_completa.csv}). Representative points that reproduce the observed anomaly are listed in Table~\ref{tab:scan}.

\begin{table}[h]
\centering
\caption{Selected points from the parameter scan that reproduce $\Delta_{\rm CKM}\simeq0.0015$.}
\label{tab:scan}
\begin{tabular}{cccc}
\toprule
$a$ [fm] & $\kappa_{\rm vac}$ & $\langle\phi\rangle$ & $V_{ud}^{\rm eff}$ \\
\midrule
0.06 & $5\times10^{-4}$ & 1.4 & 0.97352 \\
0.09 & $1\times10^{-3}$ & 0.8 & 0.97342 \\
0.06 & $5\times10^{-4}$ & 1.0 & 0.97371 \\
\bottomrule
\end{tabular}
\end{table}

A continuum extrapolation of the form $V_{ud}(a)=V_{ud}(0)+c_1a^2+\kappa_{\rm vac}\langle\phi\rangle$ yields $V_{ud}(0)\simeq0.97365$ for the Aeternvm case, while the pure Standard-Model lattice result extrapolates to $\simeq0.97420$. The deficit therefore survives $a\to0$. Figure~\ref{fig:kappa} shows the linear growth of $\Delta_{\rm CKM}$ with $\kappa_{\rm vac}$; the observed band $0.001$--$0.002$ is crossed at $\kappa_{\rm vac}=(5\pm2)\times10^{-4}$.

\begin{figure}[h]
\centering
\includegraphics[width=0.9\columnwidth]{../figures/fig1_deficit_kappa.png}
\caption{CKM deficit generated by vacuum depletion as a function of $\kappa_{\rm vac}$. The pink band marks the experimentally observed anomaly.}
\label{fig:kappa}
\end{figure}

\section{Predictions}
The framework yields three immediate, falsifiable predictions:
\begin{itemize}
\item \textbf{A-dependence:} $V_{ud}$ extracted from superallowed decays should decrease with nuclear mass number $A$ (see Paper II).
\item \textbf{Magnetic-field dependence:} $\Delta_R^V(B)=\Delta_R^V(0)+c\kappa_{\rm vac}B^2$.
\item \textbf{Neutron versus nuclear:} the existing tension between free-neutron and $0^+\to0^+$ determinations of $V_{ud}$ is of the expected size $\kappa_{\rm vac}(\langle\phi\rangle_{\rm nuc}-\langle\phi\rangle_{\rm vac})$.
\end{itemize}

\section{Conclusions}
Within the Aeternvm Vacuvm picture the CKM unitarity deficit is reinterpreted as a measurement of vacuum rigidity rather than evidence for new particles. The required coupling $\kappa_{\rm vac}\sim5\times10^{-4}$ is ultra-weak yet cumulative. Full dynamical simulations with a live $\phi$ field remain to be performed; a toy implementation is described in Paper III.

\begin{acknowledgments}
We thank the open Lattice QCD community for publicly available data and codes that made this exploratory study possible.
\end{acknowledgments}

\begin{thebibliography}{99}
\bibitem{Seng2018} C.-Y. Seng, M. Gorchtein, and M. J. Ramsey-Musolf, Phys.\ Rev.\ Lett.\ \textbf{121}, 241804 (2018).
\bibitem{Seng2020} C.-Y. Seng, X. Feng, M. Gorchtein, and L.-C. Jin, Phys.\ Rev.\ D \textbf{101}, 111301 (2020).
\bibitem{FermilabMILC} Fermilab Lattice and MILC Collaborations, arXiv:2306.XXXX (illustrative placeholder; replace with final published reference).
\bibitem{Hardy2020} J.\ C.\ Hardy and I.\ S.\ Towner, Phys.\ Rev.\ C \textbf{102}, 045501 (2020).
\bibitem{CAA} B. Capdevila et al., Phys.\ Rev.\ Lett.\ \textbf{125}, 231802 (2020).
\end{thebibliography}

\end{document}
